\subsection{Karakteristik Wirausaha di Bidang TI}
\begin{enumerate}
\item Keuntungan Wirausaha di Bidang TI:
  \begin{enumerate}
  \item Skalabilitas tinggi.
  \item Biaya operasional rendah.
  \item Pasar global.
  \item Perubahan dan inovasi yang cepat.
  \item \textit{Work From Anywhere} (WTA) atau bekerja dari mana saja.
  \item Tuntutan kebutuhan lebih tinggi.
  \item Potensi income lebih tinggi.
  \item Berpotensi menjadi pionir atau tren yang mengakibatkan dampak baru.
  \item Mudah memperoleh modal.
  \item Potensi otomatisasi.
  \end{enumerate}
\item Kelemahan Wirausaha di Bidang TI:
  \begin{enumerate}
  \item Tingkat kompetisi relatif ketat.
  \item Adopsi dan revolusi yang cepat.
  \item Biaya pengembangan dan infrastruktur cenderung tinggi.
  \item Ketergantungan pada sumber daya teknis cenderung tinggi.
  \item Risiko keamanan dan privasi.
  \item Risiko regulasi dan kepatuhan.
  \item Tingkat kegagalan startup yang tinggi.
  \end{enumerate}
\end{enumerate}

\subsection{Mengidentifikasi Peluang Usaha dalam Lingkungan Sekolah}
\begin{enumerate}
\item Aspek Penting dalam Usaha:
  \begin{enumerate}
  \item Pemahaman terhadap keahlian dam kemampuan.
  \item Observasi terhadap tren dan kebutuhan pasar.
  \item Kreativitas dalam penerapan teknologi.
  \item Pengalaman dalam proyek-proyek sekolah.
  \item Kerjasama dengan industri lokal.
  \item Pemahaman terhadap kebutuhan sekolah dan komunitas.
  \item Kesadaran terhadap permasalahan yang dihadapi masyarakat.
  \item Konsultasi dengan guru atau profesional TI.
  \end{enumerate}
\item Tahapan Menentukan Usaha:
  \begin{enumerate}
  \item Mengidentifikasi kemampuan dan keterampilan di bidang TI.
  \item Analisis tren industri TI saat ini.
  \item Pemantuauan kebutuhan pasar.
  \item Evaluasi terhadap keterbatasan dan peluang.
  \item Pembentukan ide \textit{(ideation)} dan kreativitas.
  \item Validasi konsep.
  \item Pengembangan model bisnis.
  \item Pemilihan nama dan branding.
  \item Perencanaan strategi awal.
  \end{enumerate}
\item Faktor Internal dan Eksternal:
  \begin{enumerate}
  \item Faktor Internal:
    \begin{enumerate}
    \item Sumber daya.
    \item Kapasitas produksi.
    \item Kompetensi diri.
    \item Kultur perusahaan.
    \item Visi dan misi.
    \end{enumerate}
  \item Faktor Eksternal:
    \begin{enumerate}
    \item Tren pasar.
    \item Regulasi dan kebijakan pemerintah.
    \item Teknologi.
    \item Kondisi ekonomi.
    \item Kompetisi.
    \item Faktor sosial dan struktural.
    \end{enumerate}
  \end{enumerate}
\end{enumerate}
